\conclusion

В ходе выполнения преддипломной практики
была сформулирована математическая постановка задачи
обеспечения межкластерной репликации данных
в геораспределённой системе
в условиях сетевых задержек,
отказов узлов и каналов связи,
а также возможных сетевых и инфраструктурных атак.

В работе построена формальная модель системы,
включающая описание множества кластеров,
их топологии взаимодействия,
модели журналов репликации,
случайных процессов задержек и отказов,
а также модели атакующего воздействия.
На основе введённых обозначений
сформулированы количественные критерии
целостности, доступности и согласованности данных.

Целостность данных формализована
как ограничение на вероятность нарушения корректности репликации
и обоснована применением криптографических механизмов
контроля аутентичности и неизменности записей.
Доступность описана через среднюю долю времени
успешного обслуживания операций чтения и записи,
а согласованность — через математическое ожидание
лага репликации между активным и пассивным кластерами.

Сформулирована задача выбора параметров
механизма репликации и отказоустойчивости,
минимизирующих эксплуатационные издержки
при выполнении заданных вероятностных ограничений.
Обосновано влияние настраиваемых параметров
на показатели доступности и согласованности системы,
а также показано, что требование целостности
обеспечивается на уровне протокола защиты данных.

Полученные результаты создают формальную основу
для дальнейшей разработки и экспериментальной проверки
геораспределённой системы межкластерной репликации,
устойчивой к сетевым и инфраструктурным атакам.
